\documentclass[11pt,a4paper]{article}
\usepackage{curso}
\begin{document}
\portada{Ejercicios --- Lección}{11.5}{Agentes de investigación con modelos de lenguaje}
\section{Ejercicios}
\begin{enumerate}[label=\textbf{E\arabic*.}]
\item (B) Ventajas de GRPO para las recompensas $(1;\,0{,}4;\,0;\,0{,}7)$.
\item (B) ¿Por qué GRPO no necesita red de valor?
\item (B) ¿Por qué la política solo con ejecución puso el 35\,\% en la trampa?
\item (B) ¿Qué demostraría de verdad la automejora recursiva?
\item (I) Si todos los intentos de un grupo tienen la misma recompensa, ¿qué gradiente produce GRPO? ¿Qué implica para tareas demasiado fáciles o difíciles?
\item (I) Factor de mejora de Thermo-RSI y fracción de la distancia a Opus 4.8 que se cerró.
\item (A) Ejecute \texttt{grpo\_juguete.py} con rúbrica de peso 0{,}1 y 0{,}5; ¿desde qué peso desaparece la trampa?
\end{enumerate}
\section{Soluciones}
\begin{solucion}
\textbf{E1.} Media 0{,}525, desviación 0{,}370: $(+1{,}28;\,-0{,}34;\,-1{,}42;\,+0{,}47)$.

\textbf{E2.} La media del grupo hace de línea base de la recompensa esperada.

\textbf{E3.} La trampa obtenía la ejecución máxima, igual que el método correcto: sin rúbrica no hay preferencia entre ambos.

\textbf{E4.} Que el agente descubra un algoritmo nuevo, validado de forma independiente, que mejore el hardware o el propio ciclo, y que esa mejora se repita.

\textbf{E5.} Ventajas nulas (desviación cero): no hay aprendizaje; el currículo debe tener dificultad intermedia.

\textbf{E6.} $0{,}361/0{,}127\approx2{,}8$; $(0{,}361-0{,}127)/(0{,}464-0{,}127)\approx69\,\%$.

\textbf{E7.} Verificado con 100 semillas: con peso 0 la trampa queda en 0{,}35; con 0{,}1, 0{,}3 o 0{,}5 la correcta llega a 0{,}998. En este juguete basta un peso pequeño, porque la rúbrica separa por completo la correcta de la trampa.
\end{solucion}
\end{document}
